% !TeX root = ../../Book1.tex
% 纲要标题：### 3.2 有限维性的组合后果
% 纲要位置：几何学计划纲要.md，第 211 行。
% 正文状态：初稿；尚未正式审核。

% 纲要内容开始
% 
% * 仿射依赖与 Carathéodory 定理
% * Radon 分割的线性代数证明
% * Helly 定理的归纳论证
% * 去掉维数或凸性条件后的失败
% 
% 纲要内容结束

\section{有限维性的组合后果}
\label{g1:sec:0302}

在平面上，用十个点的加权平均表示一个位置时，是否真的需要十个点？若许多凸区域每三个都有公共点，是否就能找到同时属于所有区域的点？两个问题看似一个关于表示、一个关于相交，却都受仿射维数控制。本节始终在有限维实空间中讨论；关键是把“点太多”转化为仿射依赖，再把依赖关系中的正负系数转化为凸组合。

\subsection{仿射依赖与 Carathéodory 定理}
\label{g1:subsec:0302:01}

按\secref{g1:sec:0301}的凸包定义，若 \(x\in\operatorname{conv}S\)，就有某个有限表示 \(x=\sum_{i=1}^m\lambda_i p_i\)，其中 \(p_i\in S\)、\(\lambda_i\ge0\)、\(\sum_i\lambda_i=1\)。定义没有限制 \(m\)；即使 \(S\) 很大，也要等写出具体表示之后才知道使用了多少点。Carathéodory 定理给出一个只依赖维数的统一上界。\footnote{Constantin Carathéodory（1873--1950），希腊数学家，在变分法、测度论与函数论中作出重要贡献。}

\begin{theorem}[Carathéodory 定理]\label{g1:thm:caratheodory}
	\termidx[Caratheodorydingli]{Carathéodory 定理}\foreignidx{Carathéodory's theorem}设 \(S\subseteq\mathbb R^n\) 非空，\(d=\dim\operatorname{aff}(S)\)。则每个 \(x\in\operatorname{conv}S\) 都可写成 \(S\) 中至多 \(d+1\) 个仿射无关点的凸组合。特别地，至多需要 \(n+1\) 个点。
\end{theorem}
\begin{proof}
	取 \(x=\sum_{i=1}^m\lambda_i p_i\) 的一个有限凸组合表示，删去零系数并合并重复点，使各 \(\lambda_i>0\)。如果这些点已仿射无关，证明就结束了，因为 \(d\) 维仿射空间中至多有 \(d+1\) 个仿射无关点。下面说明：只要它们仿射相关，就能减少至少一个点，同时保持表示的点 \(x\) 不变。

	回顾\cref{g1:def:affine-frame}之后的对称判别：仿射相关意味着存在不全为零的实数 \(c_i\)，满足
	\[
		\sum_{i=1}^m c_i=0,\qquad \sum_{i=1}^m c_i p_i=0.
	\]
	第二个等式是向量等式；若将所有 \(p_i\) 平移同一向量，第一个等式保证它不变。因此，把系数由 \(\lambda_i\) 改成 \(\lambda_i-tc_i\)，既不会改变系数和，也不会改变所得点。剩下的障碍只有非负性：怎样让某个系数降到零，又不让别的系数先变负？

	由于 \(\sum_i c_i=0\) 且 \((c_i)\ne0\)，其中既有正数也有负数。取
	\[
		t=\min_{c_i>0}\frac{\lambda_i}{c_i}>0,\qquad \mu_i=\lambda_i-tc_i.
	\]
	当 \(c_i>0\) 时，最小值的选择保证 \(\mu_i\ge0\)，而至少一个这样的 \(\mu_i=0\)；当 \(c_i\le0\) 时，\(\mu_i\ge\lambda_i>0\)。又有
	\[
		\sum_i\mu_i=1,\qquad \sum_i\mu_i p_i=x.
	\]
	删去零系数后，得到使用更少点的凸组合。只要剩余点仍仿射相关，就重复此操作。点数每次严格下降，故有限次后停止；此时剩余点仿射无关，最多为 \(d+1\) 个。
\end{proof}

沿仿射依赖移动全部权重，既保持点的位置和权重总和，又能用最小比值守住非负条件。这里用的是 \(S\) 的仿射维数。例如，一批点虽然写成 \(\mathbb R^{100}\) 中的坐标，但若全在同一个仿射平面内，凸包中的每个点仍只需三个点来表示。

\begin{example}[沿依赖关系减少权重]\label{g1:exm:convex-coefficient-reduction}
	令 \(p_1=(0,0),p_2=(2,0),p_3=(2,2),p_4=(0,2)\)，并取
	\[
		x=\frac18p_1+\frac38p_2+\frac14p_3+\frac14p_4=\left(\frac54,1\right).
	\]
	关系 \(p_1-p_2+p_3-p_4=0\) 的系数和为零。按证明取 \(t=\min\{1/8,1/4\}=1/8\)，便得到
	\[
		x=\frac12p_2+\frac18p_3+\frac38p_4.
	\]
	新系数非负且和为 \(1\)，故 \(x\) 落在由后三个顶点构成的三角形内。定理只保证存在一个短表示，并不保证它唯一，也不保证每次消去都会得到使用点数最少的表示。例如正方形中心可用一条对角线的两个端点表示。
\end{example}

\subsection{Radon 分割的线性代数证明}
\label{g1:subsec:0302:02}

上面的证明把仿射依赖用于调整已有权重。同一依赖式还可以直接产生两份凸组合：将它分成正、负两边，两边经同一个正数归一化，就会成为表示同一点的两个凸组合。这给出 Radon 定理。\footnote{Johann Radon（1887--1956），奥地利数学家，研究测度论、变分法和微分几何，并提出由积分数据恢复函数的 Radon 变换。}

\begin{theorem}[Radon 分割]\label{g1:thm:radon}
	\termidx[Radondingli]{Radon 定理}\foreignidx{Radon's theorem}设 \(p_1,\ldots,p_{d+2}\) 为 \(\mathbb R^d\) 中的点，允许重复。则指标集可分成两个非空、不相交的集合 \(I,J\)，使
	\[
		I\cup J=\{1,\ldots,d+2\},\qquad
		\operatorname{conv}\{p_i:i\in I\}\cap\operatorname{conv}\{p_j:j\in J\}\ne\varnothing.
	\]
	这样的指标分组称为\term[Radonfenge]{Radon 分割}[Radon partition]。
\end{theorem}
\begin{proof}
	把每个点提升为 \(\mathbb R^{d+1}\) 中的向量 \((p_i,1)\)。有 \(d+2\) 个向量而空间维数为 \(d+1\)，所以存在不全为零的实系数 \(c_i\)，使
	\[
		\sum_i c_i(p_i,1)=0.
	\]
	最后一个坐标给出 \(\sum_i c_i=0\)，其余坐标给出 \(\sum_i c_i p_i=0\)。令 \(I_+=\{i:c_i>0\}\)、\(I_-=\{i:c_i<0\}\)，它们都非空，且
	\[
		a\coloneqq\sum_{i\in I_+}c_i=-\sum_{j\in I_-}c_j>0.
	\]
	把正、负项分别移到等号两边并除以 \(a\)，得到
	\[
		z\coloneqq\sum_{i\in I_+}\frac{c_i}{a}p_i
		=\sum_{j\in I_-}\frac{-c_j}{a}p_j.
	\]
	两边系数都非负且和为 \(1\)，故 \(z\) 同时属于两组点的凸包。最后把所有零系数指标任意分配到两组，例如全部加入 \(I_+\)，令 \(I=\{i:c_i\ge0\}\)、\(J=I_-\)。加入点不会缩小凸包，因此交点 \(z\) 仍然保留，且 \(I,J\) 构成所要求的分割。
\end{proof}

这里分割的是指标。若两个指标对应同一点，仍可把它们放到不同组；这也解释了为什么定理允许重复点。提升后的向量 \((p_i,1)\) 的最后一个坐标专门用来记录系数和条件，与\secref{g1:sec:0101}中仿射组合不依赖原点的条件完全对应。

\begin{example}[平面四点的两种分割]\label{g1:exm:radon-plane}
	若四点依次为凸四边形的顶点，两条对角线相交，对角顶点各成一组，得到两段凸包的交点。若一点位于另三点所成三角形的内部，就把内部点单独作为一组，另三个点作为一组。此时交集是内部点本身，不能要求两组各有两个点。\cref{g1:fig:radon}画出了这两种情形；共线、重合等退化情形由同一线性代数证明处理，不需要另凭图形分类。
\end{example}

\begin{figure}[htbp]
	\centering
	% !TeX root = ../Book1.tex
\begin{tikzpicture}[scale=0.9]
	\begin{scope}
		\coordinate (p1) at (0,0);
		\coordinate (p2) at (3,0);
		\coordinate (p3) at (2.5,2.1);
		\coordinate (p4) at (-0.3,1.7);
		\coordinate (z) at (1.14043,0.95796);
		\draw[gray!55] (p1)--(p2)--(p3)--(p4)--cycle;
		\draw[very thick,blue!70!black] (p1)--(p3);
		\draw[very thick,orange!85!black] (p2)--(p4);
		\fill[blue!70!black] (p1) circle(2pt) node[below left] {\(p_1\)};
		\fill[blue!70!black] (p3) circle(2pt) node[above right] {\(p_3\)};
		\fill[orange!85!black] (p2) circle(2pt) node[below right] {\(p_2\)};
		\fill[orange!85!black] (p4) circle(2pt) node[above left] {\(p_4\)};
		\fill (z) circle(2pt) node[above left] {\(z\)};
		\node at (1.3,-0.7) {两条线段相交};
	\end{scope}
	\begin{scope}[xshift=6cm]
		\coordinate (q1) at (0,0);
		\coordinate (q2) at (3,0);
		\coordinate (q3) at (0.8,2.1);
		\coordinate (q4) at (1.1,0.75);
		\filldraw[draw=blue!70!black,fill=blue!8,thick] (q1)--(q2)--(q3)--cycle;
		\fill[blue!70!black] (q1) circle(2pt) node[below left] {\(p_1\)};
		\fill[blue!70!black] (q2) circle(2pt) node[below right] {\(p_2\)};
		\fill[blue!70!black] (q3) circle(2pt) node[above] {\(p_3\)};
		\fill[orange!85!black] (q4) circle(2.5pt) node[right] {\(p_4=z\)};
		\node at (1.3,-0.7) {一点位于三角形内};
	\end{scope}
\end{tikzpicture}

	\caption[平面四点的 Radon 分割]{平面四点的 Radon 分割。左图两组各有两个点，公共点 \(z\) 是对角线交点；右图一组只有 \(p_4=z\)，另一组的凸包是三角形。颜色表示分组。}
	\label{g1:fig:radon}
\end{figure}

\subsection{Helly 定理的归纳论证}
\label{g1:subsec:0302:03}

直线上的有限个非空区间若每两个相交，就有公共点。到了平面，每两个凸集相交已不够：三角形的三条边两两相交，却没有共同交点。正确的检查数量随维数增加为 \(d+1\)，这就是 Helly 定理。\footnote{Eduard Helly（1884--1943），奥地利数学家，研究凸集、积分方程与泛函分析，并在有界线性泛函的延拓理论中作出早期贡献。}

\begin{theorem}[有限族的 Helly 定理]\label{g1:thm:helly}
	\termidx[Hellydingli]{Helly 定理}\foreignidx{Helly's theorem}设 \(C_1,\ldots,C_N\) 是 \(\mathbb R^d\) 中的有限个凸集，\(N\ge1\)。若其中每个至多含 \(d+1\) 个集合的非空子族都有非空交集，则
	\[
		\bigcap_{i=1}^N C_i\ne\varnothing.
	\]
	不要求这些凸集闭、有界或有非空内部。
\end{theorem}
\begin{proof}
	当 \(N\le d+1\) 时，全族本身就是假设中的子族。第一个需要证明的情形是 \(N=d+2\)：每次删去一个集合，剩下的 \(d+1\) 个集合有公共点。我们先利用这些点和 Radon 分割找到全族的交点，再把任意较大的族压缩为少一个集合的族。

	设 \(N=d+2\)。对每个 \(i\)，选取
	\[
		p_i\in\bigcap_{j\ne i}C_j.
	\]
	点 \(p_i\) 只可能不属于 \(C_i\)。对这 \(d+2\) 个点应用\cref{g1:thm:radon}，得到分割 \(I,J\) 及两凸包的公共点 \(z\)。任意固定一个指标 \(k\)：若 \(k\in I\)，则 \(J\) 中每个指标都不等于 \(k\)，故所有 \(p_j\)（\(j\in J\)）都在 \(C_k\) 中。由 \(C_k\) 凸，整个 \(\operatorname{conv}\{p_j:j\in J\}\) 都包含在 \(C_k\) 中，特别是 \(z\in C_k\)。若 \(k\in J\)，改用 \(I\) 中各点，同样得到 \(z\in C_k\)。因此 \(z\) 属于全部 \(d+2\) 个集合。

	现在对 \(N\) 归纳。设 \(N>d+2\)，且结论对少于 \(N\) 个集合成立。把末尾两个集合合并为
	\[
		D=C_{N-1}\cap C_N,
	\]
	考虑由 \(C_1,\ldots,C_{N-2},D\) 组成的 \(N-1\) 个凸集。这样合并不改变全族的交集；需要检查的是，新族是否仍满足小子族相交条件。

	任取新族中至多 \(d+1\) 个集合。若不含 \(D\)，原假设直接给出非空交集。若含 \(D\)，把它重新展开为 \(C_{N-1},C_N\)，就得到原族中至多 \(d+2\) 个集合。展开后至多有 \(d+1\) 个时，交集由原假设非空；恰有 \(d+2\) 个时，这个子族的每个至多 \(d+1\) 元子族仍满足原假设，所以由已经证明的 \(d+2\) 情形，它的全交也非空。于是新族满足归纳假设，其全交非空，而这正是原族的全交。
\end{proof}

证明中凸性的作用很具体：所有 \(p_j\) 都在 \(C_k\) 中，才能把这些点的凸组合 \(z\) 也放进 \(C_k\)。维数则先控制产生仿射依赖所需的点数，再决定要检查的子族大小。对于 \(N\ge d+1\)，只检查恰有 \(d+1\) 个集合的子族已经足够，因为较小子族可以扩充到这个大小；定理采用“至多”的表述，是为了同时覆盖 \(N<d+1\) 的情形。

\ProseParagraphBreak

若所有集合都在一个 \(r\) 维仿射子空间中，就可通过仿射坐标把该子空间看作 \(\mathbb R^r\)，检查数量相应降为 \(r+1\)。反过来说，若有限个凸约束无共同解，必能从中挑出至多 \(d+1\) 个已经互相矛盾的约束。这一逆否表述给出了规模受维数控制的不可行子族。相关定理及进一步推广可参阅\cite[第3.2, 3.4节]{Gallier2011Geometric}。

\subsection{去掉维数或凸性条件后的失败}
\label{g1:subsec:0302:04}

这三个定理中的数量界都是最优的。先看一个同时检验这些界的模型：取 \(d\ge1\)，令 \(v_0,\ldots,v_d\) 为 \(\mathbb R^d\) 中的仿射无关点，记其单纯形为 \(\Delta\)。由\cref{g1:prop:affine-coordinates}，单纯形中的每个点都有唯一的非负重心坐标 \((\lambda_0,\ldots,\lambda_d)\)，且这些坐标之和为 \(1\)。

\begin{example}[单纯形给出的最优界]\label{g1:exm:convex-dimension-sharp}
	单纯形的重心 \(b=(v_0+\cdots+v_d)/(d+1)\) 的全部重心坐标都为正。由坐标唯一性，\(b\) 不能由少于 \(d+1\) 个顶点作凸组合得到，所以 Carathéodory 定理中的 \(d+1\) 不能统一缩小。

	同样，这 \(d+1\) 个顶点不可能有 Radon 分割：若两个不相交顶点组的凸包有公共点，把等号两边的凸组合相减，就得到一个非零、系数和为零的仿射依赖关系，与仿射无关矛盾。因此保证分割所需的 \(d+2\) 个点也不能统一减少。

	最后，对每个 \(i\) 令
	\[
		F_i=\operatorname{conv}\{v_j:j\ne i\}.
	\]
	它在重心坐标中由 \(\lambda_i=0\) 描述。任取 \(d\) 个这样的面，设未取的下标为 \(k\)，顶点 \(v_k\) 就属于所取各面的交；但全部 \(d+1\) 个面相交会要求每个 \(\lambda_i=0\)，与系数和为 \(1\) 矛盾。因此 Helly 定理不能把检查数量从 \(d+1\) 降到 \(d\)。
\end{example}

维数上界也不能彻底取消。考虑所有仅有有限个非零坐标的实数列所成的向量空间，以 \(e_1,e_2,\ldots\) 表示标准坐标向量。对任意 \(m\)，点 \((e_1+\cdots+e_m)/m\) 都属于 \(\{e_i:i\ge1\}\) 的凸包，但要从这些向量作凸组合表示它，前 \(m\) 个坐标迫使 \(e_1,\ldots,e_m\) 全部出现。因此不存在与所在有限维子空间的维数无关的统一点数界。对任意事先指定的正整数 \(k\)，再把一个 \(k\) 维单纯形的 \(k+1\) 个面放入这个空间，由上一例，它们的每个至多含 \(k\) 个集合的子族均相交，全族却不相交。因此 Helly 定理也没有与维数无关的统一检查数量。

\begin{example}[非凸集合不能只检查小子族]\label{g1:exm:helly-nonconvex}
	在直线 \(\mathbb R\) 上取
	\[
		C_1=\{0,1\},\qquad C_2=\{1,2\},\qquad C_3=\{0,2\}.
	\]
	任意两个集合的交分别为 \(\{1\},\{0\},\{2\}\)，而三个集合没有公共点。它们都是闭且有界的集合，问题出在凸性。例如 \(0,2\in C_3\)，但其中点 \(1\notin C_3\)；Helly 证明中把若干已有点换成凸组合的步骤便失去了依据。
\end{example}

\begin{example}[无限族还需要额外条件]\label{g1:exm:helly-infinite-family}
	在 \(\mathbb R\) 中令 \(C_k=[k,\infty)\)，其中 \(k=1,2,\ldots\)。每个有限子族的交都是非空闭射线，但全族的交为空，因为一个实数不能大于等于所有正整数。这些集合甚至都是闭凸集，所以“有限族”不能直接删除。

	若只补上有界性，取 \(D_k=(0,1/k)\) 仍有同样的有限相交而全交为空现象。此时每个集合都凸且有界，却在端点处不闭。有限 Helly 定理所做的是把有限次检查的数量降下来；把有限相交再推进到无限相交，是另外一个需要紧性等条件的问题。
\end{example}
